\chapter{AI agents and proposals}
\label{ch:agents}

mybibli can be opened to an AI agent --- Claude Code, Hermes Agent or
any other client of the \emph{Model Context
Protocol}\index{MCP}\index{Model Context Protocol} (MCP). The agent
reads your catalog and suggests improvements: a series and its missing
tomes, a missing description or publisher, a cover. Every suggestion
arrives as a \emph{proposal}\index{proposal} that waits on the
\textbf{Proposals} page until a librarian accepts, edits or rejects
it.

This chapter walks through:

\begin{itemize}
  \item what an agent can and cannot do,
  \item creating a \emph{propose} key and allowing the agent's
        address,
  \item connecting an MCP client --- the generic facts first, then
        Claude Code,
  \item the tools the agent sees,
  \item reviewing proposals: list, detail, comparison, edit, accept,
        reject,
  \item expected tomes and the wish list,
  \item where proposed covers wait, and how they are cleaned up.
\end{itemize}

\section{What an agent can and cannot do}
\label{sec:agents-scope}

\begin{note}
\textbf{An agent never writes the catalog.} It reads titles, series
and contributors, and it files proposals. Nothing changes in the
catalog until a Librarian or an Admin accepts a proposal on the
Proposals page. Rejecting, or simply ignoring, a proposal leaves the
catalog exactly as it was.
\end{note}

With a \emph{propose} key an agent can:

\begin{itemize}
  \item search and read titles (bibliographic fields, contributors,
        volumes with their shelf location, series positions),
        contributors and series --- including the gaps of a closed
        series;
  \item list the titles without a cover, or with incomplete metadata;
  \item propose a new series, or additions to an existing one (titles
        you hold, and \emph{expected tomes} you do not own yet);
  \item propose values for a title's subtitle, description,
        publisher, publication year, Dewey code, and contributors to
        add;
  \item propose a cover image;
  \item read the history of its own proposals, with your decisions
        and your rejection reasons.
\end{itemize}

It cannot:

\begin{itemize}
  \item create, change or delete anything in the catalog --- titles,
        volumes, series, contributors, locations, loans;
  \item see borrowers, loans, users, settings or other keys'
        proposals;
  \item erase a value: a proposal only fills or corrects a field, it
        never empties one, and it only \emph{adds} contributors;
  \item make the server download anything: a proposed cover is sent
        as an image, and the web addresses the agent cites as
        sources are shown to you as links, never fetched by mybibli.
\end{itemize}

Every proposal must cite its \textbf{sources} --- the \texttt{http}
or \texttt{https} pages where the agent found the information --- so
you can check them before deciding.

\section{Before you start: the network}
\label{sec:agents-network}

The MCP endpoint lives at \texttt{/api/mcp}\index{api/mcp@\texttt{/api/mcp}}
on the same web server as the rest of mybibli. It is meant for an
agent running on a computer \emph{of your local network}. Like the
rest of mybibli (chapter~\ref{ch:installation}), it must never be
published on the Internet: the key travels in every request, in clear
text on a plain-HTTP LAN.

\section{Step 1 --- create a \emph{propose} key}
\label{sec:agents-key}

\begin{enumerate}
  \item As an Admin, open \textbf{Administration \(\rightarrow\) API
        keys} (section~\ref{sec:api-keys-create}).
  \item Give the key a label that names the agent, e.g.\
        \emph{``Claude Code on the office laptop''}. The label is
        shown next to each proposal the key files.
  \item Choose the scope \textbf{Read + propose}\index{propose key}.
  \item Click \textbf{Create key} and copy the key from the dialog:
        it is shown only once. A propose key starts with
        \texttt{mybibli\_pr\_}.
\end{enumerate}

Only a propose key opens \texttt{/api/mcp}; a read-only or
read-write key is refused there. Use one key per agent: revoking it
(section~\ref{sec:agents-revoke}) then stops exactly that agent.

\section{Step 2 --- allow the address the agent uses}
\label{sec:agents-hosts}

To defend against \emph{DNS rebinding} (a web page tricking a browser
into talking to your NAS under a foreign name), the endpoint answers
only requests whose \texttt{Host} header is on a list. Two environment
variables, read at start-up, control it (chapter~\ref{ch:installation},
table of section~\ref{sec:env-vars}):

\begin{description}
  \item[\texttt{MYBIBLI\_MCP\_ALLOWED\_HOSTS}]\index{env vars!MYBIBLI\_MCP\_ALLOWED\_*}
        Comma-separated host names or addresses, each optionally with
        a port (\texttt{host} or \texttt{host:port}; an entry without
        a port matches any port). Default when unset or empty:
        \texttt{localhost,127.0.0.1,::1} --- enough only when the
        agent runs on the machine that hosts mybibli. A set value
        \emph{replaces} the default, so on a NAS list every name and
        address the agent uses, and keep the loopback names if a
        local client needs them:
\begin{lstlisting}
MYBIBLI_MCP_ALLOWED_HOSTS=nas.local,192.168.1.10,localhost
\end{lstlisting}
  \item[\texttt{MYBIBLI\_MCP\_ALLOWED\_ORIGINS}]\index{env vars!MYBIBLI\_MCP\_ALLOWED\_*}
        Comma-separated origins (\texttt{scheme://host:port},
        \texttt{:*} for any port). Default: empty. MCP clients such as
        Claude Code send no \texttt{Origin} header and need nothing
        here; a request that carries an \texttt{Origin} is refused
        unless it is listed. List one only for a client or proxy that
        forwards one (the MCP Inspector, section~\ref{sec:agents-inspector}).
        Browser pages cannot call the endpoint at all: it sends no
        CORS headers.
\end{description}

Set them in \texttt{.env} (the shipped \texttt{docker-compose.yml}
passes both through) and restart the container. These are deployment
settings, not Admin settings, on purpose: a page of the web interface
must not be able to widen a network defence.

\section{Step 3 --- connect the agent}
\label{sec:agents-connect}

\subsection{What any MCP client needs}
\label{sec:agents-connection-facts}

Whatever the client --- Claude Code, Hermes Agent, or another one ---
it needs these facts:

\begin{center}\small
\begin{tabular}{@{} l >{\raggedright\arraybackslash}p{9.5cm} @{}}
\toprule
\textbf{Setting} & \textbf{Value} \\
\midrule
Address   & \texttt{http://<host>:<port>/api/mcp}, where
            \texttt{<host>} is one of the allowed hosts of
            section~\ref{sec:agents-hosts} and \texttt{<port>} the
            port you open mybibli on (\texttt{HOST\_PORT}, default
            8080). \\
Transport & \emph{Streamable HTTP} (sometimes listed as just
            \emph{HTTP}). The server is stateless: every call is one
            \texttt{POST} answered with JSON; there is no event
            stream and no session to keep. \\
Header    & \texttt{Authorization: Bearer <propose key>} on every
            request. \\
Scope     & The key must be a \emph{propose} key
            (section~\ref{sec:agents-key}). No OAuth, no login. \\
Protocol  & MCP revisions \texttt{2025-11-25} and
            \texttt{2026-07-28}. Older revisions are refused. \\
\bottomrule
\end{tabular}
\end{center}

A client that can set a custom header on a streamable-HTTP server can
connect. stdio-only clients and the older ``HTTP + SSE'' transport are
not supported.

\subsection{Claude Code}
\label{sec:agents-claude-code}

Register the server once with \texttt{claude mcp add}. Replace the
address and the key with yours:

\begin{lstlisting}
claude mcp add --transport http mybibli \
  http://nas.local:8080/api/mcp \
  --header "Authorization: Bearer mybibli_pr_..."
\end{lstlisting}

The default scope (\texttt{-s local}) keeps the server --- and its
key --- in your own Claude Code configuration for the current
directory; \texttt{-s user} makes it available in every directory.
Avoid \texttt{-s project}: it writes the key into a
\texttt{.mcp.json} file meant to be shared with a repository.

Check it with \texttt{claude mcp list} (the server should show as
connected), or with \texttt{/mcp} inside a session. The tools appear
to Claude as \texttt{mcp\_\_mybibli\_\_<tool>}, e.g.\
\texttt{mcp\_\_mybibli\_\_search\_titles}.

\subsubsection*{Without registering: \texttt{--mcp-config}}

For a one-off or scripted session, describe the server in a JSON file
and pass it on the command line. Keep the file private (it holds the
key), e.g.\ \texttt{chmod 600 mybibli-mcp.json}:

\begin{lstlisting}
{
  "mcpServers": {
    "mybibli": {
      "type": "http",
      "url": "http://nas.local:8080/api/mcp",
      "headers": { "Authorization": "Bearer mybibli_pr_..." }
    }
  }
}
\end{lstlisting}

\begin{lstlisting}
claude --mcp-config mybibli-mcp.json
# non-interactive, mybibli tools only:
claude -p "Which closed series have gaps?" \
  --strict-mcp-config --mcp-config mybibli-mcp.json \
  --allowedTools "mcp__mybibli__list_series"
\end{lstlisting}

\subsection{Checking the endpoint with the MCP Inspector}
\label{sec:agents-inspector}

The MCP Inspector is the protocol's reference test client. It is the
quickest way to tell a network or key problem from an agent problem.
It runs with Node.js (\texttt{npx} downloads it on first use) on any
machine that can reach mybibli. In the commands below, replace the
address and the key.

\begin{enumerate}
  \item \textbf{List the tools} (command-line mode):
\begin{lstlisting}
npx @modelcontextprotocol/inspector --cli \
  http://nas.local:8080/api/mcp --transport http \
  --header "Authorization: Bearer mybibli_pr_..." \
  --method tools/list
\end{lstlisting}
        Expected: the twelve tools of section~\ref{sec:agents-tools}.
  \item \textbf{Call one}:
\begin{lstlisting}
npx @modelcontextprotocol/inspector --cli \
  http://nas.local:8080/api/mcp --transport http \
  --header "Authorization: Bearer mybibli_pr_..." \
  --method tools/call --tool-name get_title --tool-arg id=1
\end{lstlisting}
        Expected: the record of title~1 (or an ``unknown title''
        error if that id does not exist --- any answer from the tool
        proves the connection).
  \item \textbf{Optionally, the web interface}: run
\begin{lstlisting}
npx @modelcontextprotocol/inspector
\end{lstlisting}
        without arguments, open the address it prints, choose the transport
        \emph{Streamable HTTP}, enter the address, add the
        \texttt{Authorization} header, then \emph{Connect},
        \emph{List Tools}, and run \texttt{get\_title}. The Inspector
        talks to mybibli through its own local proxy; if that proxy
        forwards its \texttt{Origin} and the connection is refused
        with \emph{403}, add the origin it shows (typically
        \texttt{http://localhost:6274}) to
        \texttt{MYBIBLI\_MCP\_ALLOWED\_ORIGINS}, restart, and remove
        it again when you are done.
\end{enumerate}

Each tool call is logged at \texttt{info} level as an \texttt{MCP tool
call} line naming the tool and the key id (chapter~\ref{ch:operations-debugging}).
The MCP library's own messages are capped at \texttt{warn}. To see
them while you investigate, enter \texttt{info,rmcp=debug} as the
\textbf{Log level} in \textbf{Administration \(\rightarrow\) System
\(\rightarrow\) Logging} (it applies at once, no restart), then put
the previous value back. Setting \texttt{MYBIBLI\_LOG\_LEVEL} in
\texttt{.env} does not help here: the level saved in the panel wins at
every start.

\subsection{When the connection fails}

\begin{center}\small
\begin{tabular}{@{} l >{\raggedright\arraybackslash}p{10cm} @{}}
\toprule
\textbf{Answer} & \textbf{Likely cause} \\
\midrule
401 & No \texttt{Authorization} header, a mistyped key, or a revoked
      key. \\
403 \texttt{insufficient\_scope} & The key is read-only or
      read-write; create a propose key. \\
403, other & The \texttt{Host} the client uses is not in
      \texttt{MYBIBLI\_MCP\_ALLOWED\_HOSTS}, or the request carries an
      \texttt{Origin} that is not in
      \texttt{MYBIBLI\_MCP\_ALLOWED\_ORIGINS}. \\
400 / -32022 & (The second one is a JSON-RPC error,
      \emph{Unsupported protocol version}.) The client speaks a protocol revision mybibli does not serve; update the
      client. \\
406 / 415 & The client does not speak streamable HTTP (wrong
      \texttt{Accept} or \texttt{Content-Type}). \\
405 & The client tried \texttt{GET} or \texttt{DELETE} (an old
      HTTP + SSE client); choose the streamable-HTTP transport. \\
413 / 408 & A request over $\approx$\,3.7\,MiB (a cover too large) or
      longer than 30\,s. \\
\bottomrule
\end{tabular}
\end{center}

\section{The tools}
\label{sec:agents-tools}

The agent discovers the tools by itself; each carries a description
written for it. For reference:

\subsubsection*{Reading}

\begin{description}[style=nextline]
  \item[\texttt{search\_titles}] Free-text search, or an exact ISBN or
        volume label (V-code); filters by media type and genre.
  \item[\texttt{get\_title}] One title in full: fields, contributors
        and roles, volumes and their location, series positions,
        whether it has a cover, fields edited by hand.
  \item[\texttt{search\_contributors}] Contributors by name, with how
        many titles each one has.
  \item[\texttt{list\_series}] Series by name, with owned positions
        and gap count; can list only series with gaps.
  \item[\texttt{get\_series}] One series: its positions, the missing
        ones of a closed series, its expected tomes.
  \item[\texttt{list\_titles\_without\_cover}] Titles with no cover
        image.
  \item[\texttt{list\_titles\_with\_incomplete\_metadata}] Titles
        missing a description, publisher, contributor or publication
        year.
  \item[\texttt{list\_my\_proposals}] The proposals filed with
        \emph{this} key: status, decision date, your rejection reason,
        and when an equivalent proposal may be filed again.
\end{description}

\subsubsection*{Proposing}

\begin{description}[style=nextline]
  \item[\texttt{propose\_series\_create}] A new series --- name,
        closed (with its number of volumes) or open --- with the
        titles you hold at their positions and expected tomes.
  \item[\texttt{propose\_series\_attach}] Titles to place in an
        existing series, and/or expected tomes to add to it.
  \item[\texttt{propose\_title\_metadata}] Subtitle, description,
        publisher, publication year, Dewey code, and contributors to
        add, for one title.
  \item[\texttt{propose\_cover}] A cover image for one title (JPEG,
        PNG, WebP or GIF, at most 2\,MiB).
\end{description}

A tool refuses, with an explanation the agent can act on, a proposal
that could never be accepted: an unknown title, a position already
taken, a value a librarian's form would refuse, a missing source, a
second pending proposal for the same series, or a proposal you
rejected recently (section~\ref{sec:agents-delay}). Each agent may
have at most 100 pending cover proposals at a time.

\section{Reviewing proposals}
\label{sec:proposals}

\subsection{Seeing that proposals are waiting}

Librarians and Admins see a \textbf{Proposals}\index{proposals page}
entry in the navigation bar. While proposals are pending it carries a
badge with their number (\texttt{99+} above 99), and the home page's
attention section shows a \textbf{Pending proposals} indicator. Both
lead to \texttt{/proposals}. Anonymous visitors see neither.

\subsection{The list}

\texttt{/proposals} shows the pending proposals by default, newest
first, 25 per page. Two filters narrow it:

\begin{itemize}
  \item \textbf{Status} --- \emph{Pending}, \emph{Accepted},
        \emph{Rejected}, \emph{Outdated}, or all;
  \item \textbf{Kind} --- \emph{New series}, \emph{Series additions},
        \emph{Metadata}, \emph{Cover}, or all.
\end{itemize}

Each row names the subject (the proposed series name, or the series
or title concerned, linked to its page), when it was filed, by which
key, its status and its number of lines. \textbf{Open} leads to the
detail.

\subsection{The detail page}

The detail page shows who filed the proposal and when, its status,
the \textbf{Sources} the agent cited (links that open in a new tab ---
check them), and the \textbf{Proposed changes}, one card per line:

\begin{itemize}
  \item \textbf{Series proposals} --- the series line (name, type,
        number of volumes), then each title with its position (and
        the positions an omnibus covers), then each expected tome
        with its position, title and ISBN.
  \item \textbf{Metadata proposals} --- one comparison table per
        field, \textbf{Current} next to \textbf{Proposed} (for
        contributors: \textbf{To add}), side by side even on a phone.
        A badge in words says whether the field is \emph{New} (empty
        today), \emph{Changed} or \emph{Unchanged}; the proposed value
        of a new or changed field is also set off by a thick border,
        so the highlight never relies on colour alone.
  \item \textbf{Cover proposals} --- the current cover (or the
        placeholder) next to the proposed image, and a sentence that
        says whether accepting \emph{adds} a cover or \emph{replaces}
        the current one.
\end{itemize}

\subsection{Editing before you accept}

While a proposal is pending you can correct it instead of rejecting
it. \textbf{Edit} opens a line in place; \textbf{Save} checks it with
the same rules as when the agent filed it.

\begin{itemize}
  \item A new series: its name, type and number of volumes.
  \item A title line: its position, the omnibus flag and the last
        position.
  \item An expected tome: its position, title and ISBN.
  \item A metadata field: the proposed value (the current value is
        the reference and does not change); for contributors, each
        name and role --- clear a name to drop that contributor.
\end{itemize}

\textbf{Remove} takes a line out of the proposal (title, expected
tome, metadata field or contributor line). The line is not applied
when you accept, but it stays on record in a \emph{Lines removed
before the decision} section. Some lines cannot be removed: the series
line, the single line of a cover proposal, and the last line of a
metadata proposal --- reject the proposal instead. A cover cannot be
edited: a different image is a new proposal.

An edited line keeps the agent's value under it (\emph{The assistant
proposed: \dots}) and says who edited it and when. The agent's own
history still shows what it filed.

\subsection{Accepting}

\textbf{Accept} opens a confirmation that summarises what will be
written --- the series created or updated and how many titles and
expected tomes, the fields written to the title, the contributors
credited, or the cover installed. On confirmation:

\begin{itemize}
  \item a series proposal creates the series, places the titles and
        adds the expected tomes; an expected tome whose book the
        series already holds at another position is skipped and
        reported;
  \item a metadata proposal writes the fields and credits the
        contributors (one already credited in that role is skipped).
        The written fields count as \emph{edited by hand}: a later
        automatic metadata download never overwrites them;
  \item a cover proposal installs the image as the title's cover,
        also protected from automatic cover downloads.
\end{itemize}

The success message links to the series or title.

\subsection{Outdated proposals}

A proposal is checked again at the moment you accept it. If its
subject changed since the agent filed it --- the title or series was
edited or deleted, a position was taken in the meantime, a series of
the same name now exists, the cover changed, a contributor role was
deleted --- \textbf{nothing is
applied} and the proposal becomes \emph{Outdated}; the message says
why. Ask the agent for a new proposal based on the catalog as it is
now, or make the change by hand.

Accepting a proposal that writes fields or a cover changes the title,
so the other pending proposals for the same title become outdated when
you try to accept them. Decide them in the order you prefer, and ask the agent to
file again what still matters.

Revoking a key (section~\ref{sec:agents-revoke}) also marks all its
pending proposals as outdated; the detail page then says they were
decided by \emph{the system}.

\subsection{Rejecting}

\textbf{Reject} closes the proposal without changing anything. You
may give a \textbf{reason} (optional, at most 1000 characters): the
agent reads it in its history (\texttt{list\_my\_proposals}) and can
learn from it. A proposed cover is deleted when its proposal is
rejected.

\subsection{The resubmission delay}
\label{sec:agents-delay}

After a rejection, an \emph{equivalent} proposal is refused for a
while, so an agent does not file the same thing again and again.
Equivalent means the same kind and the same subject: the same series
(a new series by its name, ignoring case and accents), the same title
\emph{and} the same set of fields for metadata, the same title for a
cover. The rule applies to every key: creating a new key does not
reset it.

The delay is set in \textbf{Administration \(\rightarrow\) System
\(\rightarrow\) Proposals}\index{resubmission delay}
(section~\ref{sec:system}): from 0 (no delay) to 3650 days, default
\textbf{90 days}. A change applies to the next proposal, without a
restart.

\section{Expected tomes and the wish list}
\label{sec:expected-tomes}

An \emph{expected tome}\index{expected tome} names a tome of a series
that you do not own yet: a position, a title and, if known, an ISBN.
Expected tomes come from accepted series proposals, or you enter them
yourself.

\begin{itemize}
  \item \textbf{On the series page}, an expected tome fills its gap
        square with its title, its ISBN and an \emph{Expected}
        marker. Everyone sees it; the owned / missing counters do not
        change, since the tome is still missing.
  \item \textbf{Entering one by hand} (Librarian or Admin):
        \textbf{Add expected tome} in the series actions, or
        \textbf{Name this tome} on an empty gap. \textbf{Edit} and
        \textbf{Delete} appear under each tome.
  \item \textbf{Add to wish list} puts the tome on your wish list
        (section~\ref{sec:wishlist}) with its title, its ISBN, and a
        note naming the series and the position. Clicking twice does
        not create a second entry.
  \item \textbf{When you catalog the book} --- scan its ISBN, or
        create the title by hand with that ISBN --- mybibli recognises
        the expected tome, places the new title at that position of
        the series, and the gap disappears. If the book was on your
        wish list, the usual banner offers to remove it
        (section~\ref{sec:wishlist}). If the position is meanwhile
        held by another title, nothing is attached and a warning says
        what to check.
  \item \textbf{Expected tomes to review} --- a tome that no longer
        fits a gap (the series' number of volumes was lowered below
        its position, or a title now holds its position) is listed
        under the grid with the reason. Change its position or delete
        it.
\end{itemize}

\section{Proposed covers: the staging area}
\label{sec:agents-staging}

A proposed cover is checked like a cover you upload yourself, then
reduced to at most 400\,px wide (JPEG) --- what you review is exactly
what acceptance installs. Until you decide, it waits in a hidden
\texttt{.staging}\index{covers!.staging} folder \emph{inside the
covers directory} (\texttt{COVERS\_DIR}; in Docker,
\texttt{/app/covers/.staging} on the \texttt{mybibli\_covers} volume).
Consequences:

\begin{itemize}
  \item Nothing to configure: every install that persists its covers
        also persists the proposed ones across upgrades.
  \item A proposed image is never public: the \texttt{/covers}
        address refuses hidden folders, and the Proposals page serves
        it to librarians only, while the proposal is pending.
  \item Cleanup is automatic. Rejecting a proposal deletes its image;
        accepting moves it into place; revoking a key deletes the
        images of its pending proposals. The scheduled purge
        (section~\ref{sec:system}, \emph{auto-purge interval}) removes
        anything left behind at its next run --- the image of a
        proposal that went outdated, for instance; a filing that never
        completed only once it is an hour old.
  \item Backups of the covers volume (chapter~\ref{ch:backup}) include
        the folder; restoring it is harmless, the purge sorts it out.
\end{itemize}

\section{Revoking an agent's key}
\label{sec:agents-revoke}

In \textbf{Administration \(\rightarrow\) API keys}, revoke the key.
The agent is refused from its next request, all its pending proposals
become outdated, and their proposed covers are deleted. Its accepted
and rejected proposals stay in the history. To let the agent back in,
create a new key and update the client configuration.
